%\exer{Application: déplacement d'un cavalier sur un échiquier}
%\vspace{-0.5cm}
\section*{Exercice d'application $\star$}

\marginnote{Des exemples de graphes ou dictionnaires sont présents ici : \url{https://github.com/xpessoles/2025_2026_Enseignements/raw/main/PTSI_Informatique/TP/TP_13_ApplicationsGraphes/TP_13_ApplicationsGraphes.py}.}

On donne la carte de la région Nouvelle Aquitaine.
\begin{figure}[H]
	\centering
		\includegraphics[width=.7\linewidth]{fig_01}
	\caption{Carte et graphe de la région Nouvelle-Aquitaine}
\end{figure}


\subsection*{Dictionnaire}
On souhaite créer un dictionnaire permettant de faire une correspondance entre le numéro des départements et leur nom.

\question{Implémenter une fonction \lstinline|add_dpt(d:{}, dpt:int, nom:str) -> None| qui ajoute un département au dictionnaire  \lstinline{d}.}

\begin{exemple}[Tests]
\begin{lstlisting}
>>> dico = {}
>>> add_dpt(dico,64,"Pyrénées-Atlantiques")
>>> print(dico)
        {64:"Pyrénées-Atlantiques"}
\end{lstlisting}
\end{exemple}


\subsection*{Dictionnaire d'adjacence}
On souhaite modéliser la carte par un dictionnaire d'adjacence où les  sommets sont des numéros de départements et les valeurs sont la liste des départements limitrophes. 

\question{Donner l'instruction permettant de créer un graphe appelé \lstinline{aqui} avec les Pyrénées-Atlantiques (64) et les Landes uniquement (40).}


\begin{question}
Implémenter la fonction \lstinline|existe_d(G:{}, dpt:int) -> bool| qui renvoie \lstinline{True} si le département est dans le graphe, \lstinline{False} sinon. Tester votre fonction.
\end{question}

\question{Implémenter la fonction \lstinline|ajoute_dpt_d(G:{}, dpt:int) -> None| qui ajoute un département au graphe \lstinline{G}.  Tester votre fonction.}


\begin{question}
Implémenter la fonction \lstinline|sont_voisins_d(G:{}, dpt1:int, dpt2:int) -> bool| qui renvoie \lstinline{True} si les deux départements sont voisins, \lstinline{False} sinon.
Tester votre fonction.
\end{question}


\question{Implémenter la fonction \lstinline|ajoute_une_frontiere_d(G:{}, s1: int, s2: int) -> None| qui ajoute une arête \lstinline{G} entre deux départements limitrophes. 
Tester votre fonction.
}


\question{Implémenter la fonction \lstinline|ajoute_dpt_voisins_d(G:{}, dpt:int, voisins : [int]) -> None| qui ajoute un département au graphe \lstinline{G}. \lstinline{dpt} désigne le département à ajouter et \lstinline{voisins} la liste des départements voisins.
Tester votre fonction.}
%
%\begin{exemple}[Tests]
%\begin{lstlisting}
%>>> aqui = {}
%>>> add_dpt_d(aqui,64,[40])
%>>> print(aqui)
%        {64:[40], 40:[64]}
%>>> add_dpt_d(aqui,47,[40,33,24])
%>>> print(aqui)
%        {64:[40], 40:[64,47], 47:[40,33,24]}
%\end{lstlisting}
%\end{exemple}


\question{Implémenter la fonction \lstinline|voisins_d(G:{}, dpt:int) -> [int]| qui renvoie la liste des départements voisins de \lstinline{dpt}.
Tester votre fonction.}


\question{Implémenter la fonction \lstinline|suppr_dpt_d(G:{}, dpt:int) -> None| qui supprime un département de la région.
Tester votre fonction.}


\question{Implémenter la fonction \lstinline|frontieres_d(G:{}) -> [()]| qui renvoie la listes des arêtes du graphe. Une arête sera une tuple (couple) de deux départements limitrophes. On supprimera les doublons.
Tester votre fonction.}


\question{Implémenter la fonction \lstinline|degre_d(G:{}, dpt:int) -> int| qui donne le degré d'un sommet.
Tester votre fonction.}


\question{Implémenter la fonction \lstinline|degre_max_d(G:{}) -> int| qui donne le degré maximum du graphe.
Tester votre fonction.}


\question{Implémenter la fonction \lstinline|degre_max_l_d(G:{}) -> [int]| qui renvoie la liste des sommets de plus haut degré.
Tester votre fonction.}


\question{Implémenter la fonction \lstinline|degre_max_dpt_d(G:{}, d:{}) -> [str]| prend en argument le graphe de la région ainsi que le dictionnaire des départements qui renvoie le nom du (ou des) départements ayant le plus de départements limitrophes.Tester votre fonction.}

